%% kapa: exam=1 round=21 year=2010 period=2 form=A subject=appraisal_law subject_ko=부동산관계법규 seq=12 number=52
%% source: Q-Net 감정평가사 기출문제 (cst00309) artlSeq=5024181 file=extracted/1008/5024181/2010년 제21회 감정평가사 1차시험 기출문제/2교시.hwp
%% conversion: kordoc 4.24.51 -> tools/gen1.py (auto); review=proofread-gpt-sol; flags=none
\begin{question}{52}
국토의 계획 및 이용에 관한 법령상 도시관리계획의 입안에 관한 설명으로 옳은 것은?
\begin{choices}
  \item 특별시장․광역시장․시장 또는 군수는 인접한 관할 구역 전부 또는 일부를 포함하여 도시관리계획을 입안할 경우 입안권자를 지정하는 협의가 성립되지 아니하면 공동으로 입안하여야 한다.
  \item 국토해양부장관은 국가계획과 관련된 도시관리계획에 대하여는 관계 중앙행정기관의 장의 요청 없이도 직접 입안할 수 있다.
  \item 일정한 사항에 관해서는 주민도 도시관리계획을 입안할 수 있으며 이 경우 도시관리계획도서와 계획설명서를 첨부하여야 한다.
  \item 주민이 입안한 도시관리계획은 당해 지방자치단체에 설치된 지방도시계획위원회의 자문을 거쳐야 한다.
  \item 국토해양부장관, 시․도지사, 시장 또는 군수가 도시관리계획을 입안할 경우 주민의 의견을 청취하여야 하며, 주민의 의견 청취에 필요한 사항은 국토해양부장관이 정한다.
\end{choices}
\end{question}
\answer{2}
